\section{Discussion}

\noindent\textbf{Why decoder adaptation works.}
SAM's ViT-H encoder, pre-trained on 1B masks, already encodes texture gradients and edge patterns that are useful for COD---it simply has no decoder trained to interpret \emph{subtle} texture disruptions as boundaries. Fine-tuning teaches exactly this mapping. That a 16.6-point mIoU gain comes from training 0.6\% of parameters (4M of 636M) suggests the encoder representations are richer than the original decoder can exploit.

\noindent\textbf{Prompt robustness from mixed training.}
Alternating point and box prompts forces the decoder to learn \emph{what} the camouflaged object is, not \emph{where the user clicked}. Point-only training overfits to spatial proximity; box-only training overfits to box geometry. Mixed training breaks both failure modes, which explains why the gain transfers to prompt types never seen during training (edge, grid).

\noindent\textbf{LLPM: where it helps and why.}
The LLPM's gains concentrate on edge prompts (+2.6 mIoU) and dark/boundary-ambiguous images. The learned gate settled at $\alpha = 0.185$, a moderate perturbation that sharpens boundaries without disrupting the global image statistics the frozen encoder was trained on. Pre-encoder and decoder-side adaptation address different failure modes and combine without interference.

\noindent\textbf{Prompting context.}
SAM is an \emph{interactive} segmentation model: a user clicks on the target object. Our evaluation follows the standard SAM protocol~\cite{kirillov2023sam}, deriving prompts from ground-truth masks to simulate an idealized click. The results should be read as measuring \emph{promptable} segmentation quality, not fully automatic COD. Under this protocol, the practical question is robustness to prompt imprecision: base SAM's mIoU spans 3.5$\times$ across prompt types (0.205--0.727), while our decoder spans only 1.1$\times$ (0.694--0.776). Encoder-adaptation methods still achieve higher $S_\alpha$ and MAE (\cref{tab:comparison}), so encoder and decoder specialization address different aspects of the problem. Pairing our decoder with a lightweight automatic prompt generator (e.g., a detection head) is a natural next step toward fully automatic deployment.

\noindent\textbf{Limitations.}
(1)~Trained on COD10K and validated on CAMO and NC4K without retraining; CHAMELEON evaluation is left for future work. (2)~Requires user prompts, unlike fully automatic COD methods. (3)~Encoder-side adaptation (SAM-Adapter) reaches higher absolute performance on $S_\alpha$ and MAE. (4)~LLPM training requires live encoder forward passes (7.5 hrs vs.\ 90 min for decoder-only).
